\section{Objectives and Scope}\label{sec:intro:scope}

\subsection{Objectives}\label{sec:intro:objectives}

The work pursues objectives that operationalize the research questions in \cref{sec:intro:rq}.
First, we specify and implement a transformer-style encoder that ingests bounded windows of \glspl{pdw} from interleaved pulse streams and maps them to a dual embedding (\cref{sec:method:formulation}), together with a post-hoc clustering stage applied independently to the emitter and the mode branches.
We then study how performance depends on practical limits on context length, mixture density, and noise, reporting agreement with the ground-truth simulator partitions using standard scores such as \gls{nmi} and \gls{ari} (\cref{ch:experiments}).
Second, we formulate at least one joint training objective that couples a within-train supervised contrastive term for emitter identity with a batch-wide supervised contrastive term over the global mode labels, and compare it against a matched baseline trained primarily for deinterleaving or representation quality without that coupling, holding architecture capacity and data conditions as constant as feasible (\cref{ch:method,ch:experiments}).

\subsection{Delimitations}\label{sec:intro:delimitations}

The scope is deliberately narrow.
We operate at the level of \glspl{pdw} or equivalent tabular pulse features supplied to the model; pulse detection, \gls{rf} conditioning, antenna-level processing, and geometry-driven fusion across platforms are out of scope.
This data is generated in a propietary simiulation which different scenarios are run and pulse data is collected, inlcuding PDW and corresponing emitter and mode labels.
The focus is offline or batch analysis of recorded streams rather than hard real-time implementation on embedded hardware, and we do not address full operational emitter naming against classified libraries. Due to limitationion of the quadratic scaling of the attention mechanism used by the transformer model, we split up recorded scenarios into windows, which are evaluated individually, but in learning this is taken into account so the learning is not compleatly disjoin between windows. However the study takes this into account for possible work in the aritechtural desications, and methods to extend to online and classified libraries are possible, and will be discuessed. At inference time the trained encoder is intended to operate on previously unseen recorded scenarios whose emitter and mode memberships are both recovered by clustering the corresponding embedding trajectories, as described in \cref{sec:intro:problem}.
